% Captação (poço) e projeto/licenciamento do sistema de abastecimento
\begin{tikzpicture}[node distance=0.38cm,
    passo/.style={etapa=5.6cm, minimum height=0.85cm}]

  % --- Coluna 1: captação ----------------------------------------------
  \node[font=\footnotesize\bfseries, text=azulpsh] (h1) at (0,0) {Captação subterrânea};
  \node[passo, below=0.25cm of h1] (c1) {Estudos preliminares: viabilidade, titularidade e situação fundiária das áreas};
  \node[passo, below=of c1] (c2) {Levantamento de campo: estudos geológicos e hidrogeológicos};
  \node[passo, below=of c2] (c3) {Locação do ponto de captação e projeto construtivo do poço};
  \node[passo, below=of c3] (c4) {Anuência prévia};
  \node[passo, below=of c4] (c5) {Perfuração do poço};
  \node[passo, below=of c5] (c6) {Teste de vazão e análise físico-química da água};
  \node[passo, below=of c6] (c7) {Estudos ambientais e sociais e outorga de direito de uso da água};

  % --- Coluna 2: projeto e licenciamento ------------------------------
  \node[font=\footnotesize\bfseries, text=azulpsh] (h2) at (7.4,0) {Projeto e licenciamento};
  \node[passo] (p1) at (h2 |- c1) {Anteprojeto da solução de tratamento e armazenamento};
  \node[passo] (p2) at (h2 |- c2) {Ordenação de despesas e contratação};
  \node[passo] (p3) at (h2 |- c3) {Projeto básico: casa química, reservatório, eletromecânica e distribuição};
  \node[passo] (p4) at (h2 |- c4) {Projeto executivo};
  \node[passo] (p5) at (h2 |- c5) {Aprovação dos projetos básico e executivo};
  \node[passo] (p6) at (h2 |- c6) {Estudos ambientais e sociais para o licenciamento};
  \node[passo] (p7) at (h2 |- c7) {Obtenção das licenças ambientais};

  \coordinate (meio) at ($(c7.south)!0.5!(p7.south)$);
  \node[marco=13cm, below=0.7cm of meio] (obra) {Execução das obras de instalação (subseção~\ref*{subsec:obras})};

  \foreach \a/\b in {c1/c2,c2/c3,c3/c4,c4/c5,c5/c6,c6/c7,p1/p2,p2/p3,p3/p4,p4/p5,p5/p6,p6/p7}
    \draw[seta] (\a) -- (\b);
  % da outorga para o anteprojeto (a vazão e a qualidade da água orientam o projeto)
  \draw[seta] (c7.east) -- ++(0.5,0) |- (p1.west);
  \draw[seta] (c7.south) -- (c7.south |- obra.north);
  \draw[seta] (p7.south) -- (p7.south |- obra.north);
\end{tikzpicture}
